%! AP Statistics — Unit 1, Lesson 1.7 — Group Activity ANSWER KEY
\documentclass[10pt]{article}
\usepackage{apstats-article}
\usepackage{apstats-key}

%=============================================================================
\begin{document}

\pageheader{Unit 1, Lesson 1.7}{Group Activity --- Answer Key}
\begin{tabularx}{\linewidth}{X X X}
Name:\;\ans{Key} & Date:\; & Period:\;
\end{tabularx}

\vspace{0.10in}

\begin{tcolorbox}[colback=sky, colframe=navy, boxrule=0.8pt, arc=2mm,
    left=3mm, right=3mm, top=2mm, bottom=2mm]
\small\textbf{Dataset:} Sorted lap times (sec) for 15 swimmers:\;
62, 64, 65, 67, 68, 69, 70, 71, 72, 73, 74, 75, 77, 80, 95
\end{tcolorbox}

\vspace{0.10in}

{\large\bfseries\color{navy} Part 1 \quad Calculate Summary Statistics}

\begin{enumerate}[leftmargin=1.5em, itemsep=6pt]

  \item \textbf{Mean.}\\[3pt]
        \ans{Sum $= 62+64+65+67+68+69+70+71+72+73+74+75+77+80+95 = 1082$\\
        $\bar{x} = 1082 \div 15 \approx \mathbf{72.1}$ seconds}

  \item \textbf{Median.}\\[3pt]
        \ans{$n = 15$ (odd); the median is the 8th value: \textbf{71} seconds.}

  \item \textbf{$Q_1$, $Q_3$, IQR.}\\[3pt]
        \ans{Lower half (positions 1--7): 62, 64, 65, 67, 68, 69, 70.
        $Q_1 = \mathbf{67}$\\
        Upper half (positions 9--15): 72, 73, 74, 75, 77, 80, 95.
        $Q_3 = \mathbf{75}$\\
        $\text{IQR} = 75 - 67 = \mathbf{8}$ seconds}

  \item \textbf{Range.}\\[3pt]
        \ans{Range $= 95 - 62 = \mathbf{33}$ seconds}

\end{enumerate}

\vspace{0.10in}

{\large\bfseries\color{navy} Part 2 \quad Outlier Detection \& Choosing Statistics}

\begin{enumerate}[leftmargin=1.5em, itemsep=6pt, resume]

  \item \textbf{Fences.}\\[3pt]
        \ans{Lower fence $= 67 - 1.5(8) = 67 - 12 = \mathbf{55}$\\
        Upper fence $= 75 + 1.5(8) = 75 + 12 = \mathbf{87}$}

  \item \textbf{Is 95 an outlier?}\\[3pt]
        \ans{Yes. Since $95 > 87$ (upper fence), the value 95 seconds \textbf{is an
        outlier} by the $1.5\times\text{IQR}$ rule.}

  \item \textbf{Mean vs.\ Median; shape.}\\[3pt]
        \ans{Mean $(72.1) >$ Median $(71)$ $\Rightarrow$ the distribution is slightly
        \textbf{skewed right}. The outlier at 95 pulls the mean upward while the median
        stays near the center.}

  \item \textbf{Which pair?}\\[3pt]
        \ans{The coach should report \textbf{Median + IQR} because the distribution is
        skewed right and contains an outlier; both statistics are resistant to the
        extreme value of 95 seconds and give a more accurate picture of typical
        performance. The median of 71 seconds and IQR of 8 seconds are preferred.}

\end{enumerate}

\vspace{0.10in}

{\large\bfseries\color{navy} Part 3 \quad Individual Interpretation}

\begin{enumerate}[leftmargin=1.5em, itemsep=6pt, resume]

  \item \textbf{Mean $<$ Median implies\dots}\\[3pt]
        \ans{When mean $<$ median, the distribution is likely \textbf{skewed left} ---
        a few unusually low values pull the mean below the median.}

  \item \textbf{Which number for parents?}\\[3pt]
        \ans{The coach should report the \textbf{median (44 laps)} because the
        distribution is skewed left; the median is more resistant and better
        represents a typical swimmer's weekly performance.}

\end{enumerate}

\end{document}
